\documentclass[11pt]{article}

\usepackage[T1]{fontenc}
\usepackage[margin=0.92in]{geometry}
\usepackage{amsmath,amssymb,amsthm}
\usepackage{booktabs}
\usepackage{enumitem}
\usepackage{microtype}
\usepackage{xcolor}
\usepackage{hyperref}
\hypersetup{
  colorlinks=true,
  linkcolor=blue!55!black,
  citecolor=blue!55!black,
  urlcolor=blue!55!black,
  pdftitle={Complete EFX for Four Additive Agents and Eight Goods},
  pdfsubject={A hybrid Lean and CPC computer-assisted proof}
}
\setlength{\emergencystretch}{3em}

\newtheorem{theorem}{Theorem}[section]
\newtheorem{lemma}[theorem]{Lemma}
\newtheorem{proposition}[theorem]{Proposition}
\newtheorem{corollary}[theorem]{Corollary}
\theoremstyle{definition}
\newtheorem{definition}[theorem]{Definition}
\newtheorem{remark}[theorem]{Remark}

\newcommand{\EFX}{\mathrm{EFX}}
\newcommand{\N}{\mathcal{N}}
\newcommand{\M}{\mathcal{M}}
\newcommand{\B}{\mathcal{B}}
\newcommand{\Q}{\mathbb{Q}}

\title{\textbf{Complete EFX for Four Additive Agents and Eight Goods}\\
\large A Hybrid Lean and CPC Computer-Assisted Proof}
\author{Anonymous submission}
\date{July 2026}

\begin{document}
\maketitle

\begin{abstract}
We prove that every instance of four agents and eight indivisible goods with
nonnegative additive valuations admits a complete allocation that is
envy-free up to any good (EFX). This closes the first previously open
four-agent additive parameter cell beyond the general $m\le n+3$ frontier.

The proof is hybrid. Lean verifies a structural reduction from a hypothetical
counterexample to eight precisely specified terminal ranking cells. Each
terminal cell is translated into quantifier-free rational linear arithmetic.
An independent parser checks every SMT assertion against a separately
implemented mathematical specification, while the correspondence between
that specification and the relevant Lean predicates is an explicitly
identified human-audited trust boundary. For each formula, cvc5 emits a CPC
proof of unsatisfiability and the separate Ethos checker returns
\texttt{correct}. The eight proofs contain 13,812,082 proof-step lines and
occupy 1.668 GB. A second encoding over all 40,824 labeled onto allocations
independently closes all eight cells. An exact comparison additionally finds
that all 545,360 SMT assertions occur as the 545,360 assumptions of the CPC
proofs, with zero mismatches.

We state the complete conditional Lean interface, including a kernel-checked
perturbation argument from nonnegative to positive weights, explain the
formal-to-computational boundary, and provide a content-addressed
reproduction package. We also extract a 130-partition deletion-minimal core
for one terminal cell, with its own independently checked CPC proof, and
isolate an invariant 191-partition statement as a target for a concise human
proof. The result is a bounded computer-assisted mathematical theorem, not a
resolution of additive EFX in general or a claim that the entire proof is
internal to Lean's kernel.
\end{abstract}

\paragraph{Keywords.}
Fair division, indivisible goods, EFX, additive valuations, computer-assisted
proof, Lean, SMT, CPC.

\section{Introduction}\label{sec:introduction}

Let $\N$ be a finite set of agents and $\M$ a finite set of indivisible
goods. Agent $i$ has an additive valuation
\[
  v_i(S)=\sum_{g\in S}w_{ig},\qquad w_{ig}\ge0.
\]
A complete allocation $X=(X_i)_{i\in\N}$ partitions all goods. It is
\emph{envy-free up to any good} (EFX) if, for every pair $i\ne j$ and every
good $g\in X_j$,
\[
  v_i(X_i)\ge v_i(X_j\setminus\{g\}).
\]
Whether complete EFX allocations always exist for additive valuations is a
central open problem in discrete fair division.

Complete EFX is known for three additive agents \cite{CGM2024}. Mahara proved
complete EFX for general monotone valuations when $m\le n+3$
\cite{Mahara2021}. Hence four agents and seven goods are covered, whereas
four agents and eight goods form the first previously open four-agent
additive cell beyond that few-goods boundary. For four additive agents, an EFX
allocation leaving at most one good unallocated is known
\cite{Berger2022}, and complete EF2X is known \cite{Ashuri2025}; neither
implies complete EFX.

Our main theorem settles this fixed cell.

\begin{theorem}[Main theorem]\label{thm:main}
Every nonnegative additive valuation instance with four agents and eight
goods admits a complete EFX allocation.
\end{theorem}

The finite-search reduction is carried out for strictly positive weights:
positivity implies that every complete EFX allocation has nonempty bundles.
This does not restrict the theorem. The set of weight matrices with no
complete EFX allocation is open, because it is the complement of a finite
union of closed polyhedra, one for each complete allocation. Thus a
hypothetical counterexample with zero-valued goods remains a counterexample
after adding the same sufficiently small $\varepsilon>0$ to every item
weight. Lean verifies this closure step.

\subsection{Proof architecture}

The proof combines four layers:
\begin{enumerate}[leftmargin=2.1em]
  \item published complete-EFX inputs for four agents and seven goods and for
  three additive agents;
  \item a Lean-checked mathematical reduction to eight terminal ranking
  cells;
  \item eight independently generated and parser-audited exact
  linear-arithmetic formulas; and
  \item eight cvc5 CPC proofs accepted by the Ethos checker.
\end{enumerate}

The division of labor is deliberate. Solver output does not determine the
mathematical reduction. Conversely, Lean does not replay 13.8 million CPC
steps. The complete result is a hybrid theorem whose trusted components are
stated explicitly in Section~\ref{sec:trust}.

\subsection{Contributions}

Besides Theorem~\ref{thm:main}, the development provides:
\begin{itemize}[leftmargin=1.8em]
  \item a formal perturbation lemma extending the positive terminal proof to
  all nonnegative additive valuations;
  \item a reverse-parking criterion and a terminating threshold-cycle normal
  form for EFX allocations;
  \item a favorite-small-bundle theorem, including a constructive EFX result
  whenever the favorite map is injective and $n\le m\le2n$;
  \item formal universal-minimum and peelable residual reductions;
  \item exact canonical coverage and valid stabilizer reductions;
  \item a partition-to-perfect-matching formulation covering all 1,701
  unlabeled four-bundle partitions;
  \item independent terminal and direct-allocation encodings;
  \item exact CPC proof objects for every terminal cell; and
  \item a deletion-minimal 130-partition explanatory core.
\end{itemize}

\subsection{What the result does and does not settle}

Theorem~\ref{thm:main} is a bounded existence theorem. It does not prove EFX
for four agents and an arbitrary number of goods, nor does it resolve the
general additive EFX conjecture. Its significance is that it crosses the
first previously open four-agent parameter cell at which the existing
general few-goods theorem stops.

\section{Related work and parameter frontier}\label{sec:related}

For additive valuations, complete EFX is known for three agents
\cite{CGM2024}. For any number of agents, Mahara proved existence for
$m\le n+3$ under general monotone valuations \cite{Mahara2021}. The
four-agent partial result of Berger et al.\ leaves at most one good
unallocated \cite{Berger2022}. Ashuri, Gkatzelis, and Sgouritsa prove EF2X
for four agents and explicitly distinguish that relaxation from the open EFX
problem \cite{Ashuri2025}.

Neoh and Teh study counting and variants of EFX, recording the same
$m\le n+3$ few-goods frontier \cite{NeohTeh2025}. Br\^anzei reports strong
empirical performance of a local-search heuristic on additive instances but
does not obtain unrestricted existence \cite{Branzei2025}. Recent positive
results for hypergraph valuations and for two types of goods remain
structured subclasses of the general additive setting
\cite{Lianeas2026,Davidiuk2026}.

Akrami et al.\ give nonexistence examples from $m\ge n+5$ for monotone and
submodular valuations and formally verify their SAT foundations in Lean
\cite{Akrami2026}. They also prove complete EFX existence for three agents
and seven goods under monotone valuations. That positive result is likewise
an $m=n+4$ cell and uses certified finite search, so the broad proof strategy
has direct recent precedent. The present result is its four-agent additive
analogue, using exact linear arithmetic and CPC/Ethos rather than
SAT/DRAT. Their counterexamples are nonadditive and begin one good later.
The current landscape is therefore
\[
\begin{array}{c}
\text{general monotone existence through }m=n+3,\\[2mm]
\text{monotone existence at }(3,7),\\[2mm]
\boxed{\text{nonnegative additive existence at }(n,m)=(4,8)},\\[2mm]
\text{nonadditive failure from }m=n+5.
\end{array}
\]
Akrami et al.\ explicitly leave the $m=n+4$ and four-or-more-agent additive
frontiers open. Our theorem closes one fixed subcase of those questions; it
is not a general solution and does not claim methodological priority.

\section{Model and elementary structure}\label{sec:model}

\begin{definition}[Complete allocation]
An allocation is a function assigning every good to exactly one agent. Its
bundles $X_i$ are pairwise disjoint and have union $\M$.
\end{definition}

\begin{definition}[EFX]
A complete allocation $X$ is EFX if
\[
  v_i(X_i)\ge v_i(X_j\setminus\{g\})
\]
for every distinct pair $i,j$ and every $g\in X_j$.
\end{definition}

For an agent $i$ and bundle $S$, define the one-good-removal threshold
\[
  \tau_i(S)=\max_{g\in S}v_i(S\setminus\{g\}),\qquad
  \tau_i(\varnothing)=0.
\]
For additive valuations,
\[
  \tau_i(S)=v_i(S)-\min_{g\in S}v_i(g)
\]
when $S$ is nonempty. An allocation is EFX exactly when
\[
  v_i(X_i)\ge\tau_i(X_j)\quad\text{for all }i,j.
\]

\subsection{Reverse parking}

Suppose $X$ is EFX and $f\in X_h$. Removing $f$ only makes the owner
comparison for $h$ harder; every other target bundle either stays fixed or
shrinks. Consequently, $f$ is removable from $h$ precisely when
\[
  v_h(X_h)-v_h(f)\ge\max_{j\ne h}\tau_h(X_j).
\]
Choosing a least-valued good in $X_h$ gives:

\begin{lemma}[Reverse parking]\label{lem:parking}
An owner $h$ has a removable good if and only if
\[
  \tau_h(X_h)\ge\max_{j\ne h}\tau_h(X_j).
\]
\end{lemma}

If no owner satisfies this inequality, choose for every owner $h$ a strict
threshold-maximizing target $c(h)$. The resulting functional graph contains
a directed cycle. Rotating bundles along that cycle preserves EFX and
strictly increases the finite potential
\[
  \Phi(X)=\sum_h\tau_h(X_h).
\]

\begin{theorem}[Threshold-cycle normal form]\label{thm:cycle}
Every complete EFX allocation with nonempty bundles can be transformed by
EFX-preserving bundle rotations into one having an owner-removable good.
\end{theorem}

The formal theorem is generic over finite agents and goods and nonnegative
additive valuations. In the positive $4\times8$ setting, an empty bundle
would force every other bundle to have cardinality at most one, contradicting
completeness; hence all bundles are nonempty.

\subsection{From nonnegative to positive weights}

Let $\mathcal A$ be the finite set of complete allocations. For a fixed
$X\in\mathcal A$, the set of weights for which $X$ is EFX is an intersection
of finitely many closed linear half-spaces. Hence
\[
  \{w:\text{$w$ admits a complete EFX allocation}\}
  =
  \bigcup_{X\in\mathcal A}\{w:\text{$X$ is EFX}\}
\]
is closed, and the set of counterexamples is open.

\begin{lemma}[Positive approximation]\label{lem:nonnegative}
If a nonnegative additive $4\times8$ instance has no complete EFX allocation,
then some strictly positive additive $4\times8$ instance has no complete EFX
allocation.
\end{lemma}

\begin{proof}
For $\varepsilon\in\mathbb R$, set
$w^\varepsilon_{ig}=w_{ig}+\varepsilon$. The path
$\varepsilon\mapsto w^\varepsilon$ is continuous and passes through $w$ at
$\varepsilon=0$. Openness of the no-EFX set gives a radius $\delta>0$ such
that $w^\varepsilon$ still has no EFX allocation whenever
$|\varepsilon|<\delta$. Taking $0<\varepsilon<\delta$ makes every coordinate
strictly positive.
\end{proof}

Lean proves the closedness of each fixed-allocation EFX set, the finite-union
closedness of the existence set, continuity of the uniform perturbation, and
Lemma~\ref{lem:nonnegative}. The remainder of the reduction may therefore
work with positive weights without narrowing Theorem~\ref{thm:main}.

\subsection{Normalization and generic strictness}

Multiplying one agent's valuation row by a positive scalar preserves every
EFX comparison. We normalize each row so that
\[
  \sum_{g\in\M}w_{ig}=1.
\]
There are finitely many complete allocations. If a normalized instance has
no EFX allocation, choose a strict failed inequality for each allocation.
These finitely many failures persist in an open neighborhood. A sufficiently
small perturbation within the row-sum hyperplanes preserves positivity and
nonexistence while avoiding both coordinate-equality hyperplanes within each
row and equality of any two full valuation rows.

\begin{lemma}[Generic normalized counterexample]\label{lem:perturb}
If a strictly positive additive $4\times8$ instance has no EFX allocation,
then there is a row-normalized counterexample in which every row has pairwise
distinct item values and the four valuation rows are pairwise distinct.
\end{lemma}

Lean proves the simultaneous perturbation, rather than using genericity as an
informal without-loss-of-generality step.

\subsection{Favorite small bundles}

\begin{lemma}[Favorite-small-bundle lemma]\label{lem:favorites}
Let $X$ be a complete allocation in a positive additive instance. If every
agent owns one of its favorite goods and every bundle has size at most two,
then $X$ is EFX.
\end{lemma}

\begin{proof}
After removing any good from a target bundle, at most one good remains. The
observer values that good no more than its own favorite, which is contained
in its own bundle.
\end{proof}

The formal result is constructive. If the favorite map is injective and
$|\M|\le2|\N|$, it assigns every agent its favorite and distributes at most
one residual good per agent. Thus the injective-favorite case is closed at
$4\times8$.

\section{Reduction to eight terminal cells}\label{sec:reduction}

Assume for contradiction that a nonnegative additive $4\times8$ instance has
no complete EFX allocation. First apply Lemma~\ref{lem:nonnegative}, then
normalize and perturb the resulting positive instance using
Lemma~\ref{lem:perturb}.

\subsection{Residual cuts from smaller theorems}

Two regions reduce to published smaller cases.

\paragraph{Universal minimum.}
If some good is minimum for every agent, delete it and apply Mahara's
complete $4\times7$ theorem. The threshold-cycle normal form moves the
smaller allocation to an appropriate unenvied source; inserting the deleted
minimum good there preserves every EFX comparison.

\paragraph{Peelable agent-good pair.}
If an agent-good pair satisfies the formal peelability inequalities, delete
that agent and that good, apply the complete three-agent theorem to the
remaining seven goods, and reinsert the pair.

The peelable predicate excludes the deleted agent as an observer in the
smaller instance but retains every active-observer comparison required after
reinsertion. This observer scope was corrected during adversarial
specification review and is now part of the checked formal statement.

\begin{proposition}[Residual exclusion]\label{prop:residual}
A normalized generic counterexample has neither a universally minimum good
nor a peelable agent-good pair.
\end{proposition}

\subsection{Favorite multiplicities}

By Lemma~\ref{lem:favorites}, the four favorite goods cannot be distinct.
With strict rankings, every noninjective favorite multiplicity is, up to
agent and good relabeling,
\[
  2+1+1,\qquad3+1,\qquad2+2,\qquad4.
\]
Canonical favorite vectors for the first three are
\[
  (0,0,1,2),\qquad(0,0,0,1),\qquad(0,0,1,1).
\]
If all four agents favor good $0$, refine by the multiplicities of their
second favorites:
\[
  1+1+1+1,\qquad4,\qquad3+1,\qquad2+2,\qquad2+1+1.
\]
This produces exactly eight terminal ranking cells.

\subsection{Valid stabilizers}

Once a ranking prefix is fixed, only symmetries preserving that prefix
remain valid. Agents with identical prescribed prefixes and goods absent from
every prefix can be sorted lexicographically. Lean proves that every orbit
has a representative satisfying the corresponding stabilizer inequalities.
No symmetry exchanging inequivalent favorite or second-favorite labels is
used.

\begin{theorem}[Eight-cell reduction]\label{thm:eightcells}
Assume the standard four-agent/seven-good and three-agent additive EFX
inputs. Every positive additive $4\times8$ counterexample yields a normalized
valuation matrix satisfying one of the eight terminal cell specifications in
Table~\ref{tab:cells}.
\end{theorem}

\begin{table}[ht]
\centering
\small
\begin{tabular}{@{}lll@{}}
\toprule
Cell & Ranking data & Stabilizer\\
\midrule
favorite-$211$ & $(0,0,1,2)$ & none\\
favorite-$31$ & $(0,0,0,1)$ & favorite preserving\\
favorite-$22$ & $(0,0,1,1)$ & favorite preserving\\
second-$1111$ & $(1,2,3,4)$ & prefix preserving\\
second-$4$ & $(1,1,1,1)$ & prefix preserving\\
second-$31$ & $(1,1,1,2)$ & prefix preserving\\
second-$22$ & $(1,1,2,2)$ & prefix preserving\\
second-$211$ & $(1,1,2,3)$ & prefix preserving\\
\bottomrule
\end{tabular}
\caption{The eight terminal ranking cells. The last five share favorite
vector $(0,0,0,0)$.}
\label{tab:cells}
\end{table}

\section{Partitions and perfect matchings}\label{sec:matching}

There are exactly
\[
  S(8,4)=1,701
\]
unlabeled partitions of eight goods into four nonempty bundles. The Lean
development verifies the restricted-growth enumeration and proves that every
onto labeled allocation is obtained from one enumerated partition by a
permutation of the bundles.

Fix a partition $\B=(B_0,B_1,B_2,B_3)$. Construct a bipartite graph between
agents and bundles. Put edge $iB_b$ when
\[
  v_i(B_b)\ge\max_{c\ne b}\tau_i(B_c).
\]

\begin{proposition}[Matching equivalence]\label{prop:matching}
The partition $\B$ supports a complete EFX allocation if and only if its
agent-bundle feasibility graph has a perfect matching.
\end{proposition}

\begin{proof}
A perfect matching assigns each bundle to one distinct agent. The edge
condition says that the assigned owner passes every one-good-removal
comparison against the three other bundles, hence the allocation is EFX.
Conversely, the owner assignment of any complete EFX allocation gives
exactly such a matching.
\end{proof}

Therefore a counterexample must obstruct all $4!=24$ matchings for every one
of the 1,701 partitions. This is a finite disjunction of linear comparisons,
which can be encoded in exact rational linear arithmetic.

\section{Exact terminal formulas}\label{sec:formulas}

For each cell $c$, let $F_c$ be a quantifier-free rational
linear-arithmetic formula over the 32 item-value variables. It contains:
\begin{itemize}[leftmargin=1.8em]
  \item 32 strict positivity assertions and four row-normalization equations;
  \item pairwise row inequalities and the exact favorite or second-favorite
  order;
  \item cell-preserving stabilizer inequalities;
  \item exclusions of universal-minimum and peelable residual regions;
  \item 27,216 agent-bundle feasibility-edge definitions; and
  \item 40,824 perfect-matching obstruction assertions.
\end{itemize}

\begin{proposition}[Semantic implication]\label{prop:semantic}
For every terminal cell $c$, any valuation matrix satisfying the Lean
terminal specification for $c$ satisfies $F_c$.
\end{proposition}

This proposition describes the load-bearing bridge. The current assurance is
hybrid:
\begin{enumerate}[leftmargin=2.1em]
  \item a formula generator independent of the original Z3 encoder emits
  $F_c$;
  \item a separate parser reconstructs and checks every SMT assertion against
  a hand-written mathematical specification;
  \item the parser independently enumerates all 1,701 partitions and all 24
  matching clauses per partition;
  \item canonical vectors are checked directly against their Lean source
  constants; and
  \item the remaining correspondence between the hand-written specification
  and Lean's residual predicates, stabilizers, feasible edge, and matching
  obstruction has been manually reviewed.
\end{enumerate}

Thus Proposition~\ref{prop:semantic} is not presently a single reflected Lean
theorem. It is an explicit trusted software-and-review boundary. This is
stronger than relying on an unchecked generator, but weaker than
kernel-checked extraction from Lean.

\section{CPC certificate closure}\label{sec:certificates}

For each formula $F_c$, cvc5 emits a CPC proof and Ethos checks the proof
under pinned CPC signatures. Table~\ref{tab:proofs} gives the final proof
sizes.

\begin{table}[ht]
\centering
\small
\begin{tabular}{@{}lrr@{}}
\toprule
Cell & Proof bytes & Proof-step lines\\
\midrule
favorite-$211$ & 187,295,060 & 1,890,346\\
favorite-$31$ & 459,896,059 & 3,327,798\\
favorite-$22$ & 338,399,888 & 2,702,364\\
second-$1111$ & 21,220,067 & 249,283\\
second-$4$ & 449,335,513 & 3,374,635\\
second-$31$ & 115,319,771 & 1,186,530\\
second-$22$ & 78,871,768 & 872,213\\
second-$211$ & 18,119,773 & 208,913\\
\midrule
Total & 1,668,457,899 & 13,812,082\\
\bottomrule
\end{tabular}
\caption{CPC proof sizes. Ethos returns \texttt{correct} for every row.}
\label{tab:proofs}
\end{table}

\subsection{Formula-to-proof assumption audit}

A valid proof for a different inconsistent formula would not establish
unsatisfiability of $F_c$. To exclude this stale-proof attack, the CPC term
definitions were expanded and every proof assumption was compared with the
SMT assertion stream. Across all eight cells:
\[
  545,360\text{ formula assertions}
  =
  545,360\text{ proof assumptions},
\]
with zero mismatches. The release package includes the executable comparison
and a manifest-bound JSON result.

\subsection{Direct-allocation cross-check}

As non-load-bearing corroboration, a second formulation directly negates
every labeled onto allocation. There are
\[
  4!\,S(8,4)=40,824
\]
such allocations. This formulation does not use the partition feasibility
graph. All eight direct formulas are UNSAT: six close in Z3, while the two
hardest close in cvc5 after an audited identity-only SMT-LIB syntax
normalization.

The direct results increase confidence that no partition-to-matching coding
error drives the main certificate, but the six Z3 results do not carry
independently checked proof objects and are not part of the load-bearing
deduction.

\section{Formal composition}\label{sec:formal}

Lean proves the structural theorem through the terminal cells. The complete
exported interface is:

\noindent\begin{minipage}{\linewidth}
\footnotesize
\begin{verbatim}
theorem admitsEFX_4x8_of_standard_published_inputs_nonnegative
    (four_agent_seven_good :
      FourAgentSevenGoodStandardInput)
    (three_agent_seven_good :
      ThreeAgentSevenGoodStandardInput)
    (terminal_cells_unsat :
      EightTerminalCellsUnsat)
    (weight : Weight48)
    (nonnegative : forall agent good, 0 <= weight agent good) :
    AdmitsEFX weight
\end{verbatim}
\end{minipage}

Here \texttt{EightTerminalCellsUnsat} is the proposition that each of the eight
terminal specifications is empty. It is discharged externally by
Proposition~\ref{prop:semantic} and the CPC/Ethos results, not by a theorem
internal to Lean.

The axiom audit of the conditional Lean theorem reports exactly
\[
  [\texttt{propext},\ \texttt{Classical.choice},\ \texttt{Quot.sound}],
\]
with no \texttt{sorryAx} and no native-computation axiom in the dependency
closure.

\begin{theorem}[Hybrid closure]\label{thm:hybrid}
Assume:
\begin{enumerate}[label=(\roman*),leftmargin=2.2em]
  \item the two published smaller-case EFX theorems;
  \item the soundness of Proposition~\ref{prop:semantic}; and
  \item correctness of Ethos under the pinned CPC signatures.
\end{enumerate}
Then Theorem~\ref{thm:main} follows from the eight archived CPC proofs.
\end{theorem}

\begin{proof}
The CPC proofs and Ethos verdicts establish $\neg F_c$ for every terminal
cell $c$. By Proposition~\ref{prop:semantic}, no Lean terminal specification
is inhabited, giving \texttt{EightTerminalCellsUnsat}. Applying the
conditional Lean theorem with the two published inputs yields
\texttt{AdmitsEFX weight} for every
nonnegative $4\times8$ weight matrix. Internally, the exported theorem first
uses Lemma~\ref{lem:nonnegative} to reduce a hypothetical nonnegative
counterexample to a positive one.
\end{proof}

\section{The 130-partition explanatory core}\label{sec:core}

The favorite-$211$ cell has favorite vector $(0,0,1,2)$. Starting from all
1,701 partition obstructions, assumption-core extraction produced 200
partitions. Deletion shrinking then produced a 130-partition core.

\begin{table}[ht]
\centering
\begin{tabular}{@{}lrr@{}}
\toprule
Bundle-size shape & Initial core & Final core\\
\midrule
$3+2+2+1$ & 98 & 59\\
$2+2+2+2$ & 51 & 20\\
$4+2+1+1$ & 30 & 30\\
$3+3+1+1$ & 20 & 20\\
$5+1+1+1$ & 1 & 1\\
\midrule
Total & 200 & 130\\
\bottomrule
\end{tabular}
\caption{Shape counts in the explanatory core.}
\label{tab:core}
\end{table}

Every single-partition deletion is satisfiable, so every survivor is
indispensable relative to this core. This is deletion minimality, not
minimum-cardinality optimality.

The final core has two pieces:
\begin{enumerate}[leftmargin=2.1em]
  \item 110 partitions in which shared favorite good $0$ is a singleton,
  favorite goods $0,1,2$ are separated, and no nonfavorite good is a
  singleton; and
  \item 20 partitions of shape $2+2+2+2$ in which good $0$ is paired with one
  nonfavorite good.
\end{enumerate}

The reduced formula has a 220,777,714-byte CPC proof containing 2,032,963
proof-step lines, accepted by Ethos.

These 130 partitions lie in a simpler invariant family $\mathcal A$:
\begin{enumerate}[leftmargin=2.1em]
  \item all 131 partitions in which good $0$ is a singleton, goods
  $0,1,2$ are separated, and no nonfavorite good is a singleton; and
  \item all 60 partitions into four pairs in which goods $0,1,2$ are
  separated.
\end{enumerate}
Thus $|\mathcal A|=191$. The certificate proves the following weaker but
more symmetric statement, because its 130 obstructions are a subset of
$\mathcal A$.

\begin{proposition}[Anchored contested-favorite lemma]
\label{prop:anchored}
Let $w$ be a positive, row-normalized $4\times8$ valuation with distinct
rows, strict favorite vector $(0,0,1,2)$, no universally minimum good, and no
peelable agent-good pair. Then some partition in $\mathcal A$ has a feasible
perfect matching.
\end{proposition}

Proposition~\ref{prop:anchored} is currently a certificate corollary, not a
human proof. The Lean development defines its exact invariant family and
proves that this single statement suffices to eliminate the favorite-$211$
terminal cell. It is the most concrete analytic replacement target produced
by the core analysis.

Exact deletion witnesses have been extracted for the first 34 of the 130
partitions.
For each completed deletion, the removed partition is the sole supported
partition within the 130-core instance. It need not be unique in the full
1,701-partition universe. Successful removed partitions also typically have
several perfect matchings.

The Hall data also rejects a tempting compression. In the resulting 4,386
obstructed feasibility graphs there are 426 distinct graph signatures.
Although the full four-agent Hall witness is the most frequent primary
witness, an exact set-cover calculation needs ten different nontrivial agent
masks to cover the sample, and that size-ten cover is unique. Frequencies of
three dominant masks therefore do not support a three-case proof. The
promising compression is geometric---prove
Proposition~\ref{prop:anchored} using the two bundle families and threshold
inequalities---rather than a taxonomy of Hall masks.

\section{Trust boundary and limitations}\label{sec:trust}

The complete deduction has the form
\[
\begin{aligned}
&\text{two pinned published EFX theorems}\\
&+\ \text{Lean-checked conditional structural reduction}\\
&+\ \text{parser and human audited Lean-to-SMT implication}\\
&+\ \text{eight CPC proofs accepted by Ethos}\\
&\Longrightarrow\ \text{nonnegative additive EFX at }(4,8).
\end{aligned}
\]

\subsection{Trusted components}

The remaining trusted components are:
\begin{itemize}[leftmargin=1.8em]
  \item the mathematical proofs in the two published input papers;
  \item the Lean kernel, pinned toolchain, mathlib, and pinned dependencies;
  \item the Ethos checker and pinned CPC signatures;
  \item the parser implementation; and
  \item the manually reviewed correspondence between the parser's
  mathematical specification and the relevant Lean definitions.
\end{itemize}

\subsection{Limitations}

\begin{itemize}[leftmargin=1.8em]
  \item The entire proof is not one theorem checked inside Lean's kernel.
  \item The two published paper proofs are not re-formalized.
  \item The theorem treats only four agents and eight goods.
  \item The explanatory core does not replace the full terminal
  certification.
  \item Same-host reproduction does not substitute for independent-machine
  reproduction.
  \item Novelty has been searched in current primary literature but requires
  ordinary expert and peer-review confirmation.
\end{itemize}

\subsection{Independent internal adversarial reviews}

Two read-only internal reviewers worked independently. The mathematical
review found no material reduction error, conditional on terminal UNSAT. The
formal/computational review found no formula or proof mismatch and contributed
the 545,360-assumption comparison. Both reviewers independently identified
the semantic bridge as the highest-risk boundary. These reviews are evidence
of adversarial inspection, not external peer review.

\section{Reproducibility and artifact package}\label{sec:artifacts}

The final content-addressed package contains:
\begin{itemize}[leftmargin=1.8em]
  \item a deterministic source archive containing the Lean project,
  independent generators, and executable audits;
  \item the eight exact terminal formulas and CPC proofs;
  \item pinned CPC signatures and tool acquisition hashes;
  \item the semantic, published-input, direct-allocation, and
  formula-to-proof audits;
  \item the Lean build and axiom-audit transcript;
  \item the two internal referee reports;
  \item this submission manuscript; and
  \item machine-readable manifests and replay instructions.
\end{itemize}

Every load-bearing entry is hash-bound by the end-to-end manifest. The source
archive and manifest are identified by SHA-256 digests in the release notes.
The archived CPC proofs retain their original two absolute signature include
lines. The release verifier creates a temporary relocated copy by rewriting
only those two lines, verifies that every remaining proof-body byte is
unchanged, checks all 51 signature files against the official cvc5 1.3.3
source archive, and runs Ethos on the relocated copy. The full eight-cell
relocation replay passed on the release host, preserving every proof-body byte;
this is verifier-level relocatability, not an assumption about the archived
paths or an independent-host replication.

\paragraph{Minimal verification.}
From the certification root:
\begin{verbatim}
shasum -a 256 -c EFX48_END_TO_END_MANIFEST.sha256
\end{verbatim}
The release README gives platform-specific commands for the Lean build,
semantic audit, formula-to-proof comparison, and Ethos checks.

\section{Conclusion}\label{sec:conclusion}

Every nonnegative additive instance with four agents and eight goods admits a
complete EFX allocation. The proof closes the first previously open
four-agent cell beyond the $m\le n+3$ existence frontier. It is a fixed-cell
result, not a solution for four agents with arbitrary goods. Its mathematical
reduction is formalized in Lean, while its terminal arithmetic is discharged
by exact CPC proofs checked by Ethos under an explicit hybrid trust boundary.

The remaining mathematical challenge is explanatory rather than existential:
give a human proof of Proposition~\ref{prop:anchored}, then use the same
bundle geometry to compress the other seven cells. The remaining validation
challenge is independent reproduction and external expert review.

\appendix

\section{Terminal formula schema}\label{app:schema}

For a fixed cell, the formula uses rational variables $w_{ig}$ for
$i\in\{0,1,2,3\}$ and $g\in\{0,\ldots,7\}$.

\paragraph{Base valuation constraints.}
\[
  w_{ig}>0,\qquad \sum_{g=0}^{7}w_{ig}=1,
  \qquad \bigvee_{g=0}^{7}(w_{ig}\ne w_{jg})\quad(i\ne j).
\]

\paragraph{Residual exclusions.}
For every good, at least one agent ranks another good below it, excluding a
universal minimum. For every agent-good pair, the negation of the formal
peelability conjunction is asserted.

\paragraph{Ranking cells.}
The favorite or favorite-and-second-favorite inequalities encode one row of
Table~\ref{tab:cells}. The strict ranking inequalities make the prescribed
favorites unique.

\paragraph{Feasible edges.}
For each partition $\B$, observer $i$, and candidate own bundle $B_b$, an
edge variable is definitionally equivalent to
\[
  \bigwedge_{\substack{c\ne b\\g\in B_c}}
  \left(
    \sum_{h\in B_b}w_{ih}
    \ge
    \sum_{h\in B_c\setminus\{g\}}w_{ih}
  \right).
\]

\paragraph{Matching obstruction.}
For every $\pi\in S_4$:
\[
  \bigvee_{i=0}^{3}\neg\,\mathrm{edge}(i,\pi(i)).
\]
The conjunction over all 24 permutations says that the partition has no
perfect matching.

\section{Published-input interface}\label{app:published}

The two external inputs are:
\begin{enumerate}[leftmargin=2.1em]
  \item complete EFX for four agents and seven goods, obtained from Mahara's
  $m\le n+3$ theorem; and
  \item complete EFX for three additive agents and seven goods, obtained from
  the three-agent theorem of Chaudhury, Garg, and Mehlhorn.
\end{enumerate}

The pinned audits check exact PDF hashes, theorem fragments, completeness,
and the EFX convention. They do not verify the papers' proofs. The Lean
lifting transports the standard statements to the subtype-indexed deleted
instances used by the residual cuts.

\section{Reviewer checklist}\label{app:review}

A mathematical reviewer should check:
\begin{enumerate}[leftmargin=2.1em]
  \item the EFX convention and complete-allocation definition;
  \item the uniform perturbation from nonnegative to positive weights;
  \item row normalization and simultaneous generic perturbation;
  \item the favorite-small-bundle theorem;
  \item universal-minimum and peelable reinsertion;
  \item canonical favorite and second-favorite coverage;
  \item stabilizer soundness;
  \item partition-to-perfect-matching equivalence; and
  \item the exact two published theorem variants.
\end{enumerate}

A formal/computational reviewer should check:
\begin{enumerate}[leftmargin=2.1em]
  \item the Lean root imports and axiom audit;
  \item the parser's complete SMT assertion coverage;
  \item the manually reviewed Python-specification-to-Lean correspondence;
  \item formula and proof hashes;
  \item the 545,360-assumption comparison;
  \item CPC signature relocation;
  \item Ethos verdicts; and
  \item replay on an independent machine.
\end{enumerate}

\begin{thebibliography}{99}

\bibitem{CGM2024}
B. R. Chaudhury, J. Garg, and K. Mehlhorn.
\newblock EFX exists for three agents.
\newblock \emph{Journal of the ACM}, 71(1), 2024.
\newblock \url{https://arxiv.org/abs/2002.05119}.

\bibitem{Mahara2021}
R. Mahara.
\newblock Extension of additive valuations to general valuations on the
existence of EFX.
\newblock In \emph{ESA 2021}, LIPIcs 204, Article 66, 2021.
\newblock \url{https://doi.org/10.4230/LIPIcs.ESA.2021.66}.

\bibitem{Berger2022}
B. Berger, A. Cohen, M. Feldman, and A. Fiat.
\newblock Almost full EFX exists for four agents.
\newblock In \emph{AAAI 2022}, pages 4826--4833.
\newblock \url{https://doi.org/10.1609/aaai.v36i5.20410}.

\bibitem{Ashuri2025}
A. Ashuri, V. Gkatzelis, and A. Sgouritsa.
\newblock EF2X exists for four agents.
\newblock In \emph{AAAI 2025}, pages 13555--13563.
\newblock \url{https://doi.org/10.1609/aaai.v39i13.33480}.

\bibitem{NeohTeh2025}
J. Q. Neoh and N. Teh.
\newblock Understanding EFX allocations: Counting and variants.
\newblock In \emph{AAAI 2025}, 2025.
\newblock \url{https://doi.org/10.1609/aaai.v39i13.33536}.

\bibitem{Branzei2025}
S. Br\^anzei.
\newblock Computing envy-free up to any good (EFX) allocations via local
search.
\newblock arXiv:2510.05429, 2025.
\newblock \url{https://arxiv.org/abs/2510.05429}.

\bibitem{Akrami2026}
H. Akrami, A. Mayorov, K. Mehlhorn, S. Srinivas, and C. Weidenbach.
\newblock A counterexample to EFX: $n\ge3$ agents, $m\ge n+5$ items,
submodular valuations via SAT-solving.
\newblock arXiv:2604.18216, 2026.
\newblock \url{https://arxiv.org/abs/2604.18216}.

\bibitem{Lianeas2026}
T. Lianeas, A. Sgouritsa, and M. M. Sotiriou.
\newblock EFX allocation in (multi)hypergraphs.
\newblock In \emph{AAAI 2026}, 2026.
\newblock \url{https://doi.org/10.1609/aaai.v40i20.38759}.

\bibitem{Davidiuk2026}
V. Davidiuk, Y. Dementiev, A. Ignatiev, and D. Sagunov.
\newblock EFX and PO allocation exists for two types of goods.
\newblock In \emph{AAAI 2026}, 2026.
\newblock \url{https://doi.org/10.1609/aaai.v40i20.38723}.

\end{thebibliography}

\end{document}
